\documentclass[11pt]{article}
\usepackage{amsmath,amssymb}
\title{Original normalization and an explicit boundary transfer test}\author{}\date{}
\begin{document}\maketitle
\paragraph{Boundary bookkeeping input.}
Use the sf boundary insertion and wilson normalization calculation: its raw partition ratio, moving-minimum
second insertion, multiweight cross terms and distinction between
$\Gamma_2'$ and $m_2-m_1^2$. The boundary operator fixed
entirely by prescribed spatial links is a live scalar, not a fluctuating
interior insertion; its source derivatives remain.
Primary sources and their limits are as cited there. The deductions
below make the normalization and one failure of uniform transfer explicit.

\paragraph{1. Bulk subtraction in the original normalized Wilson convention.}
The lattice action in Luescher--Weisz, hep-lat/9504006, equation5.1,
uses a sum over both ordered orientations $\mu,\nu$. With our sum over
geometric plaquettes and $s_p=1-\operatorname{ReTr}U_p/N$, this is
\[
 S= B\sum_p s_p,\qquad B=\frac{2N}{g_0^2}.                    \tag{1}
\]
Write the original fixed-boundary expansion as
\[
 \Gamma=B V_R+F_{1,R}+B^{-1}C_R+\cdots .                      \tag{2}
\]
In the conventional expansion
$\Gamma=g_0^{-2}\Gamma_0+\Gamma_1+g_0^2\Gamma_2+\cdots$,
\[
 \Gamma_0=2N V_R,\qquad\Gamma_1=F_{1,R},\qquad
                         \Gamma_2=C_R/(2N).                 \tag{3}
\]
If the SAME-scheme bare coupling is
$g_0^2=g^2+z_1g^4+z_2g^6+\cdots$, elementary inversion gives
\[
 g_0^{-2}=g^{-2}-z_1+(z_1^2-z_2)g^2+\cdots .                 \tag{4}
\]
Thus the putative two-loop bulk-renormalized coefficient in our units is
\[
 \boxed{C_R^{\rm bulk,ren}
           =C_R+(2N)^2(z_1^2-z_2)V_R.}                       \tag{5}
\]
For a flat harmonic boundary tangent $h$, $V_R''[h,h]=E(h)$, so its
Hessian subtraction is $(2N)^2(z_1^2-z_2)E(h)$; for SU(2) the
conversion factor is $16$. No value/sign for the logarithmic $z_i$ is
imported from another coupling scheme here. Background ghost and
quantum-field subrenormalizations and boundary insertions still belong
to the completed coefficient. Equation(5) is the required algebraic
normalization, NOT proof of its finite-relative-box boundedness.
The canonical quantum response retains the overall factor $B^{-2}$.

\paragraph{2. A literal harmonic test on a single cube face.}
Take $R\ge4$ and prescribe one unit color tangent on the tangential
edge
\[
 e=((0,1,1,1),\mu=1),
\]
with every other boundary edge zero. Give the free interior initially
zero values. Exactly five original plaquettes meet $e$: two in plane $12$,
two in plane $13$, and one inward plaquette in plane $01$. The outward $01$
plaquette lies outside the cube and is absent. The four $12/13$ plaquettes
lie entirely in the fixed $x_0=0$ face. Each carries unit squared curl.
Consequently the zero-interior extension has energy5 and the FIXED
face energy is4.

Let $h$ be the genuine original Maxwell-harmonic extension of this
boundary datum, on the original relative gauge quotient. Minimization
cannot increase energy, and cannot change the four fixed curls. Hence
\[
 4\le E(h)\le5,\qquad
 \boxed{\frac{E_{\rm face}(h)}{E(h)}\ge\frac45.}               \tag{6}
\]
After normalizing total energy to one this is an admissible harmonic
direction with nonvanishing boundary curvature uniformly in $R$.
This is a purely finite original-incidence argument; no asymptotic
harmonic extension or continuum boundary trace estimate is used.

The spatial SF operator has half the boundary plaquette allocation,
so the analogous ratio has its declared factor $1/2$, still uniformly
positive. With the usual quadratic action $E/2$, its corresponding
Hessian is likewise the literal face form with that allocation.
Thus the electric/magnetic face improvement terms cannot all be
replaced by an $O(1/R)$ operator error on the entire unit-energy
boundary tangent class. The smooth fixed-profile SF expansion is not
contradicted: this test varies on the lattice boundary scale.
Boundedness is not disproved; only uniform SMALLNESS is.

\paragraph{3. The exact remaining transfer requirement.}
The original unmodified Wilson coefficient must retain both(5)'s bulk
kinetic subtraction and the actual boundary-insertion ledger if it is
compared to an improved SF coefficient. The latter includes
$c_1\partial_c\Gamma_1+c_2\partial_c\Gamma_0+
\frac12c_1^2\partial_c^2\Gamma_0$, and all source derivatives.
For all-relative cubes, the change of boundary/ghost domains and
corner terms is a separate unproved comparison; neither a periodic
bulk theorem nor a two-face slab computation supplies it.
The raw source-zero scalar cannot be reconstructed from a running
coupling derivative. No uniform estimate of the complete packet,
its scale sum, the original nonlinear remainder or S94 follows.
\end{document}
